\section{Methodology}\label{sec:method}

To rigorously evaluate the performance characteristics of on-device deep learning inference across cross-platform and native stacks, we design a full-factorial, controlled empirical benchmark following pre-registered hypotheses and established systems evaluation standards~\cite{georges2007statistically,barrett2017virtual}.

\subsection{Experimental Design and Factors}
Our benchmark evaluates five controlled independent experimental factors, summarized in Table~\ref{tab:factors}:
\begin{enumerate}
    \item \textbf{Platform ($P$):} Flutter (Dart Ahead-Of-Time compiled) and Android Native (Kotlin compiled with R8 full-mode shrinking).
    \item \textbf{Model ($M$):} Three standard edge deep learning architectures spanning computer vision and natural language processing: MobileNetV3-Large~\cite{howard2019searching}, EfficientDet-Lite0~\cite{tan2020efficientdet}, and MobileBERT-QA~\cite{sun2020mobilebert}.
    \item \textbf{Precision ($Q$):} 32-bit floating-point (FP32) and full 8-bit integer post-training quantization (INT8)~\cite{jacob2018quantization}.
    \item \textbf{Hardware Delegate ($D$):} CPU multi-threading (4 threads with XNNPACK kernel acceleration), mobile GPU (OpenCL/OpenGL compute shaders), and Android Neural Networks API (NNAPI).
    \item \textbf{Device Tier ($T$):} Physical handsets spanning entry-level, mid-range, and flagship hardware profiles.
\end{enumerate}

\begin{table}[t]
\caption{Full-Factorial Experimental Design Matrix}
\label{tab:factors}
\centering
\begin{tabular}{ll}
\toprule
\textbf{Factor} & \textbf{Evaluated Levels} \\
\midrule
Platform & Flutter (Dart 3.12.2), Android Native (Kotlin 1.9.24) \\
Models & MobileNetV3-Large, EfficientDet-Lite0, MobileBERT-QA \\
Tasks & Image Classification, Object Detection, SQuAD 1.1 QA \\
Precision & FP32 (single precision), INT8 (quantized) \\
Delegates & CPU (4 threads, XNNPACK), GPU, NNAPI \\
Sessions & 8 independent sessions per configuration \\
Iterations & 20 warm-up + 100 measured per session \\
Canonical Inputs & 50 standardized images, 50 standardized text pairs \\
\bottomrule
\end{tabular}
\end{table}

\subsection{Benchmark Applications and Parity Contract}
Ensuring strict parity between cross-platform and native code is essential to guarantee that measured differentials reflect framework overhead rather than implementation divergence. Both benchmark applications strictly adhere to the \textit{Harness Contract}:
\begin{itemize}
    \item \textbf{Identical Runtime Engine:} Both applications link against Google LiteRT 2.16.1 (ADR-0002, ADR-0008). Flutter binds to the native shared C library via Dart FFI (`tflite\_flutter`), while Android Native binds via JNI.
    \item \textbf{Cryptographic Model Verification:} Model flatbuffer files (`*.tflite`) are pre-registered with cryptographic SHA-256 digests in `models/manifest.yaml`. At application initialization, both apps compute the runtime digest of the asset binary and abort execution if any hash discrepancy exists.
    \item \textbf{Standardized Input Sequence:} Both applications bundle identical asset sets (50 vision JPEGs and 50 text question/context pairs), processed in identical deterministic sequential order.
    \item \textbf{Realistic Developer Path:} In adherence to real-world software engineering practices, Flutter performs input image decoding and bilinear resizing using the pure Dart `image` package, representing typical Flutter application architecture. Android Native uses Android framework bitmap decoding.
    \item \textbf{Timed Region Isolation:} In accordance with contract specifications, no logging, disk I/O, garbage collection triggers, or UI allocations occur within the measured pipeline timestamps $t_0 \dots t_3$. Measurement records are buffered in heap memory and flushed to disk asynchronously after configuration completion.
    \item \textbf{Release Optimization:} All measurements are obtained exclusively from release-mode binaries (`flutter build apk --release` and Android R8-optimized release APKs). Debug builds (which employ JIT compilation) are strictly excluded.
\end{itemize}

\subsection{Pipeline Stage Definitions and Metric Formalism}
To isolate the precise origin of cross-platform overhead, we instrument each iteration across four discrete timestamps:
\begin{equation}
t_0 \xrightarrow{\text{preprocess}} t_1 \xrightarrow{\text{invoke}} t_2 \xrightarrow{\text{postprocess}} t_3
\end{equation}
where $t_0$ marks the availability of the raw asset in memory, $t_1$ marks input tensor buffer completion, $t_2$ marks the completion of the runtime `run()` invocation, and $t_3$ marks the production of the typed result object.

The individual latency metrics are formally defined as:
\begin{align}
\text{preprocess\_ms} &= \frac{t_1 - t_0}{10^6} \quad (\text{ms}) \\
\text{inference\_ms} &= \frac{t_2 - t_1}{10^6} \quad (\text{ms}) \\
\text{postprocess\_ms} &= \frac{t_3 - t_2}{10^6} \quad (\text{ms}) \\
\text{end\_to\_end\_ms} &= \frac{t_3 - t_0}{10^6} \quad (\text{ms})
\end{align}

\textbf{Boundary Overhead Measurement (ADR-0006):} The total invocation latency ($\text{inference\_ms}$) perceived by application code wraps input buffer copying, foreign function call transitions, kernel execution, and output buffer extraction. To isolate the pure kernel inference time, we interrogate the runtime's internal invoke timer:
\begin{equation}
\text{inference\_native\_ms} = \begin{cases}
\frac{\text{lastNativeInferenceDurationMicroSeconds}}{1000} & \text{(Flutter)} \\[4pt]
\frac{\text{lastNativeInferenceDurationNanoseconds}}{10^6} & \text{(Native)}
\end{cases}
\end{equation}
We derive the language-boundary marshalling overhead directly as:
\begin{equation}
\text{boundary\_overhead\_ms} = \text{inference\_ms} - \text{inference\_native\_ms}
\end{equation}
A timer residual check $\text{timer\_residual\_ms} = \text{end\_to\_end\_ms} - (\text{pre} + \text{inf} + \text{post})$ is monitored to verify timestamp continuity.

\textbf{System Resource Metrics:} Total memory consumption is recorded as Proportional Set Size ($\text{pss\_mb}$) via Android's `ActivityManager.getProcessMemoryInfo()`. Process CPU utilization ($\text{cpu\_pct}$) is computed via user and system execution time deltas in `/proc/self/stat`. Thermal throttling states ($\text{thermal\_status}$) are tracked via `PowerManager.getCurrentThermalStatus()`.

\begin{table*}[t]
\caption{Model Registry and Cryptographic Integrity Specifications}
\label{tab:models}
\centering
\begin{tabular}{lllllll}
\toprule
\textbf{Model ID} & \textbf{Task} & \textbf{Input Shape} & \textbf{Precision} & \textbf{Size (MB)} & \textbf{SHA-256 Checksum} & \textbf{Source Reference} \\
\midrule
mobilenetv3\_large & Classification & $[1, 224, 224, 3]$ & FP32 & 21.89 & \texttt{f924c190e9112c7db6630ab522c8fb2972be408180d2d72b8038e1a2340b2266} & Howard et al.~\cite{howard2019searching} \\
mobilenetv3\_large & Classification & $[1, 224, 224, 3]$ & INT8 & 5.67 & \texttt{1d76f35adb879b2cde9fe39000ede3d36e6d317027799c6d0927570a22c1e8d1} & Jacob et al.~\cite{jacob2018quantization} \\
efficientdet\_lite0 & Detection & $[1, 320, 320, 3]$ & FP32 & 1.62 & \texttt{70ae8c44c4e00b62c53557405e40c719632a503dbd49703df2bf30da96bf43ab} & Tan et al.~\cite{tan2020efficientdet} \\
efficientdet\_lite0 & Detection & $[1, 320, 320, 3]$ & INT8 & 0.42 & \texttt{70c5d02357c568d512d76125eef7d687f58a1a9514c301fd25b22675fc910c95} & Tan et al.~\cite{tan2020efficientdet} \\
mobilebert\_qa & Question Ans. & $[1, 384]$ & FP32 & 20.09 & \texttt{10358b8795d14a18c5fc4ba1eebe11fd4331109b155cca9bf013877d1b2cfa60} & Sun et al.~\cite{sun2020mobilebert} \\
mobilebert\_qa & Question Ans. & $[1, 384]$ & INT8 & 5.06 & \texttt{79471b95588d440ee959f19b275667187dc134ee2be426c30f34ea02972360f7} & Sun et al.~\cite{sun2020mobilebert} \\
\bottomrule
\end{tabular}
\end{table*}

\subsection{Experimental Protocol and Environmental Controls}
Measurement sessions follow a strict execution regimen to eliminate thermal and operating system noise:
\begin{enumerate}
    \item \textbf{Device Conditioning:} Devices are charged between 50\% and 90\%, switched to Airplane Mode with Wi-Fi, Bluetooth, and GPS disabled. Screen brightness is fixed at 50\%, and display sleep is prevented via wake locks.
    \item \textbf{Randomized Execution Order:} For each of the 8 sessions, configuration ordering is shuffled using a deterministic linear congruential generator initialized with seed $S = 20261109 + \text{session}$, eliminating systematic sequence bias.
    \item \textbf{Thermal Gating and Mandatory Cooldown:} Before initiating each configuration, the runner verifies that $\text{thermal\_status} \le \text{light}$. Following each configuration, a mandatory 60-second cooldown is enforced. Any session encountering $\text{thermal\_status} \ge \text{severe}$ is quarantined and excluded.
    \item \textbf{Warm-Up Discard:} Each configuration executes 20 initial warm-up iterations to stabilize dynamic voltage and frequency scaling (DVFS) and instruction caches, followed by 100 timed measurement iterations.
\end{enumerate}

\subsection{Statistical Unit of Analysis and Hypothesis Testing}
A common error in mobile systems benchmarking is treating consecutive iterations within a single process execution as independent replicates, which inflates sample size and produces artificially deflated $p$-values (pseudo-replication).

Following ADR-0007, our primary unit of analysis is the \textbf{session median}:
\begin{equation}
\tilde{X}_{c, s} = \operatorname{median}_{i=1 \dots 100} \{ x_{c, s, i} \}
\end{equation}
yielding $n_F = 8$ and $n_N = 8$ independent replicates per configuration $c$ for Flutter and Native Android respectively. For cold-start latency, each process invocation represents an independent launch and serves as the analysis unit.

For each configuration and metric, we assess significance using the two-sided Mann--Whitney $U$ test. To quantify effect magnitude without parametric assumptions, we compute Cliff's delta ($\delta$)~\cite{cliff1993dominance}:
\begin{equation}
\delta = \frac{\#(\tilde{X}_F > \tilde{Y}_N) - \#(\tilde{X}_F < \tilde{Y}_N)}{n_F \cdot n_N}
\end{equation}
accompanied by the median ratio $\mathcal{R} = \operatorname{median}(\tilde{X}_F) / \operatorname{median}(\tilde{Y}_N)$ with 10,000-resample bootstrap 95\% confidence intervals.

Family-wise error rates are controlled across multiple comparisons using the Holm--Bonferroni procedure~\cite{holm1979simple} applied within each research question family at $\alpha = 0.05$. We establish a pre-registered criterion for \textit{practical significance}:
\begin{equation}
|\delta| \ge 0.33 \quad \text{AND} \quad \text{CI}_{95\%}(\mathcal{R}) \cap [0.95, 1.05] = \emptyset
\end{equation}
requiring at least a medium effect size and a bootstrap ratio entirely outside the $\pm 5\%$ equivalence band.
